% 347_Gell_Mann_Nishijima_Galois_En_ch.tex — auto-generated stub
% Full EN version to follow; De version is authoritative
\chapter*{Gell-Mann--Nishijima from the Galois Structure}
\addcontentsline{toc}{chapter}{Gell-Mann-Nishijima from Galois}
\begin{center}\large\textbf{Johann Pascher}\\\normalsize ORCID 0009-0000-6518-4064\\4~September 2026\end{center}
\providecommand{\BM}{\textbf{[B]}}\providecommand{\KM}{\textbf{[K]}}\providecommand{\SM}{\textbf{[S]}}
\section*{Abstract}The Gell-Mann--Nishijima relation $Q=I_3+Y/2$ is derived from two established Galois properties: (1)~Hypercharge $Y=-1+\tfrac{4}{3}\mathrm{NQR}_{13}$ from the Legendre symbol mod~13 (Doc.~346, Theorem~A). (2)~Isospin $I_3=\pm\tfrac{1}{2}$ from the Su/Sd split via the jE7 projector (Doc.~344, Theorem~E). The combination gives all four fermion charges $\{0,-1,+\tfrac{2}{3},-\tfrac{1}{3}\}$ exactly. Theorem~E [S] of Doc.~346 is closed: \textbf{[B]}. See German version for full proofs.

\section{What remains open: SU$(2)_L$ and CKM \SM}

\textbf{SU$(2)_L$ chirality.}
Chirality requires $\gamma_5 = e_1\cdots e_6 \in \text{Cl}(6)$.
The $\mathbb{Z}_2$ parity in GF$(27)^*$ separates particles from anti-particles
($k<13$ vs.\ $k\geq13$), not left-handed from right-handed.
All four primitive orbits $O_1\ldots O_4$ are same-chiral.
What is needed: the bridge GF$(27)^*\to\text{Cl}(6)\xrightarrow{\gamma_5}\text{chirality}$.
The embedding of Galois orbits into the Clifford grade structure is begun
in Docs.~341/344 but not yet complete. \SM

\textbf{CKM mixing angles.}
GF$(27)$ carries a trace bilinear form
$\langle a,b\rangle = \mathrm{Tr}_{\mathrm{GF}(27)/\mathrm{GF}(3)}(ab) = ab+(ab)^3+(ab)^9$.
The Cabibbo angle between generations would be the overlap
$|\langle O_{2,\text{Gen.1}},\,O_{4,\text{Gen.2}}\rangle|/(\|O_2\|\|O_4\|)$.
Experimental: $\sin\theta_C\approx0.225$; Galois candidate $\sin(\pi/14)\approx0.222$ (1.2\,\%).
Full derivation requires: (1) explicit trace form computation for all orbit pairs;
(2) generation assignment within orbits; (3) CKM matrix evaluation.
A concrete, executable next step requiring its own document. \SM

\begin{thebibliography}{9}
\bibitem{P344} J.~Pascher, \textit{Doc.~344}, FFGFT corpus (Sept.\ 2026).
\bibitem{P346} J.~Pascher, \textit{Doc.~346}, FFGFT corpus (Sept.\ 2026).
\end{thebibliography}
